% Fragment: regularity, branching, fusion-ring reduction, and characters.
% The main document supplies theorem environments, reviewnote, and artifact.

\section{Regularity from a positive conformal-vector identity}
\label{sec:regularity}

This section concerns the actual PCT-stable subalgebra $U\subset B$ defined
in the construction. Its conclusions do not require properness. We use
the convention $Y(a,z)=\sum_{n\in\mathbb Z}a_nz^{-n-1}$ and write
$L^e(n)=e_{n+1}$ for the Virasoro operators of a conformal vector $e$.
The subalgebras used below need not contain the ambient conformal vector.
We call an inclusion $A\subset V$ \emph{grading compatible} if the
$A$-weight of each homogeneous vector of $A$ equals its $V$-weight.
Thus a mode $a_n$ changes both gradings by $\operatorname{wt}(a)-n-1$;
the equality of the two energy operators is not required on all of $V$.
All subalgebras in the positive identity below are grading compatible
and stable under the ambient PCT involution.

\begin{certificate}[Embedded structural data]
\label{cert:structural-inputs}
The following are assertions about actual fields and their exact modes
in the lattice construction, in addition to the finite closure ledger.
\begin{enumerate}
\item The Potts fields lie in $U$, are PCT stable, and generate the
regular unitary algebra $K(\mathfrak{sl}_2,3)$, with its translation
and intrinsic grading inherited from the ambient VOA. The exact
records are \artifact{hi_cft_20260923_proof_potts.json} and
\artifact{hi_cft_20260923_proof_potts_exact.json}.
\item The actual fields $z,a$ of weights two and three lie in $U$,
are PCT stable, have central charge $8/5$, and satisfy the stress,
primary, and translation identities in
\artifact{hi_cft_20260923_proof_pf13_exact.json},
\artifact{hi_cft_20260923_proof_pf13_membership.json}, and
\artifact{hi_cft_20260923_proof_remaining_coset.json}.
Their full singular first rows agree with the parafermion rows at
$c=8/5$, $\lambda=14/473$, as polynomials in $z,a,a_1a,a_1(a_1a)$.
The polynomial records are
\artifact{hi_cft_20260923_proof_pf13_upper_poles.json},
\artifact{hi_cft_20260923_proof_pf13_nonlinear.json},
\artifact{hi_cft_20260923_proof_pf13_truncation.json},
\artifact{hi_cft_20260923_proof_pf13_simple_pole_algebra.json}, and
\artifact{hi_cft_20260923_proof_pf13_dly.json}.
The membership record also gives the positive identity
\eqref{eq:positive-stress}.
\item The lattice--Potts core $T$ below is embedded with its stated
charge realization. The six joint spectral vectors used to construct
its order-thirteen extension are actual vectors of $U_2$, with the
eigenvalues printed below. Their exact source is
\artifact{hi_cft_20260923_next_structure_core_extension.json}; the
construction of $T$ is described in
\artifact{hi_cft_20260923_proof_report.md}, Section~4.
\item The zero-mode action on the complementary eighty currents is
irreducible, and the four mixed lattice tops of dimension eighty-one
are irreducible as $U$-tops. These are exact
field-action assertions, as recorded in
\artifact{hi_cft_20260923_proof_current_module_v2.json} and
\artifact{hi_cft_20260923_next_81top_all.json}, respectively.
\end{enumerate}
All polynomial equalities here mean exact vector equalities in the
fixed lattice realization, not projected or floating-point matches.
\end{certificate}

The input from the local construction is the identity
\begin{equation}
 \omega_U=\omega_{A_{1,1}}
   +\frac{3403}{4563}\omega_{A_{1,39}}
   +\frac{25}{6}\overline p+\frac{25}{26}\overline z.
 \label{eq:positive-stress}
\end{equation}
Here $p$ and $z$ are the conformal vectors of actual embedded copies of
the three-state Potts algebra $K(\mathfrak{sl}_2,3)$ and the parafermion
algebra $K(\mathfrak{sl}_2,13)$; a bar denotes normalized averaging over
the compact current group. Each subalgebra is regular and unitary, and
its intrinsic grading agrees with the ambient grading on its states.
The averages take place in the finite-dimensional space $U_2$.
More concretely, $p,z$ lie in the already constructed current-stable
space $S_2$, whose invariant subspace is spanned by the four mutually
orthogonal stresses $\omega_{A_{1,1}},\omega_{A_{1,39}},e,f$.
For either vector $v$, its Haar average is therefore
\[
 \overline v=\sum_{w\in\{\omega_{A_{1,1}},\omega_{A_{1,39}},e,f\}}
       \frac{\langle v,w\rangle_g}{\langle w,w\rangle_g}\,w.
\]
Thus the exact coefficients in \eqref{eq:positive-stress} can be
verified inside $S_2$ without assuming $S_2=U_2$ or properness.
The compact orbit of a conformal vector has compact convex hull in the
real space of PCT-fixed weight-two vectors. Its Haar average belongs
to that hull. Carath\'eodory's theorem therefore replaces each average
by a finite convex combination of orbit points. In particular,
\begin{equation}
 \omega_U=\sum_{i=1}^N c_i e_i,\qquad c_i>0.
 \label{eq:finite-positive-stress}
\end{equation}
Zero coefficients, if any, are omitted. All simple modules of each
regular subalgebra $A_i$ have nonnegative lowest weights, with zero
occurring only for its vacuum module.

\begin{reviewnote}{R1: the embedded regular factors are local inputs}
Certificate~\ref{cert:structural-inputs} states the model-specific
field inputs explicitly. Proposition~\ref{prop:pf13-finite-reconstruction}
below proves the transport from the complete checked parafermion rows
to the actual factor identification. It does not replace verification
of the listed lattice identities, Potts embedding, or positive stress
identity. These inputs are separate from the finite closure theorem.
\end{reviewnote}

\begin{proposition}[Finite reconstruction of the parafermion factor]
\label{prop:pf13-finite-reconstruction}
Let $V,V'$ be simple vertex algebras of CFT type generated by a
conformal vector $L$ and a weight-three primary $a$. Set
$b=a_1a$ and $d=a_1b$, of weights four and five. Suppose the full
singular products $a_na,a_nb,a_nd$, $n\geq0$, agree as homogeneous
differential normal polynomials in $L,a,b,d$, with the same normal
ordering. Suppose also that their self-row is
\[
 a_5a=\frac c3\mathbf1,\quad a_4a=0,\quad a_3a=2L,\quad
 a_2a=TL,\quad a_1a=b,\quad a_0a=\frac12Tb-\frac1{12}T^3L,
\]
and that $a_nb=0$ for $n\geq4$, $a_3b\in\mathbb Ca$, and
$a_2b\in\mathbb CTa$. Then there is an isomorphism $V\cong V'$
identifying $L,a,b,d$. In particular $a_1d$ and both simple poles
must be specified polynomials; they cannot be left as independent
new generators.
\end{proposition}
\begin{proof}
We first recover the remaining generator OPEs without a circular
normal-product reduction. The Virasoro rows follow by commuting
Virasoro modes through $b=a_1a$ and $d=a_1b$:
\[
 L(m)b=(2m+1)a_{m+1}a,\qquad
 L(m)d=(2m+1)a_{m+1}b+a_1L(m)b\quad(m\geq1).
\]
These are linear in generators and derivatives under the displayed
hypotheses. Thus Virasoro creation operators can be moved to the
outside of any normal word using already known commutators.
Sesquilinearity removes derivatives. Since three non-Virasoro
factors have weight at least nine, the possible two-factor cores
through weight eight are exactly
\[
\begin{array}{c|c}
\text{weight bound}&\text{cores besides individual fields}\\\hline
6&:aa:\\
7&:aa:,\ :ab:\\
8&:aa:,\ :ab:,\ :bb:,\ :ad:.
\end{array}
\]
Reversing the order of a core uses skew symmetry and its already
known pair bracket. Commuting an external $a_n$ or $b_n$ through
outer Virasoro creation operators reduces their number; any resulting
negative mode is itself a normal product.

Write $A_r$ for the evaluation of $a_nP$, $n\geq0$, on every normal
polynomial $P$ of weight at most $r$, and $B_r$ for the corresponding
assertion with $b_n$. We use the noncommutative Wick identity \cite[Section~2]{DSK}
\[
 [x_\lambda:yz:]=:[x_\lambda y]z:+:y[x_\lambda z]:
              +\int_0^\lambda[[x_\lambda y]_\mu z]\,d\mu.
\]
The following order is acyclic. First $A_6$ is known: the only
two-factor core is $:aa:$, and $[a_\lambda a]$ is linear in the
vacuum and derivatives of $L,b$, whose brackets with $a$ are given.
The iterate identity
\begin{equation}
 b_ns=a_1(a_ns)-a_0(a_{n+1}s)
          +a_{n+1}(a_0s)-a_n(a_1s)
 \label{eq:pf13-b-iterate}
\end{equation}
with $s=b$ now determines $b_nb$, because each inner $a_jb$ has
weight at most six. Next $A_7$ follows: its new core $:ab:$ uses
only the just-computed $b,b$ row in the Wick integral.
Equation~\eqref{eq:pf13-b-iterate} with $s=d$ then determines
$b_nd$, since every inner $a_jd$ has weight at most seven.

At this point $B_6$ is known. For its core $:aa:$ the Wick integral
brackets a coefficient of $[b_\lambda a]$, of weight at most six,
with $a$, which is $A_6$ after skew symmetry. Single-field inputs
use the known $b,b$ and $b,d$ rows. We obtain $A_8$ next: its core
$:bb:$ uses $B_6$ on coefficients of $[a_\lambda b]$, and its core
$:ad:$ uses the known $L,d$ and $b,d$ rows. Finally $B_7$ follows
from $B_6$ on the coefficients of $[b_\lambda a]$ in the new core
$:ab:$. With $s=d$, the second iterate identity
\[
 d_ns=a_1(b_ns)-a_0(b_{n+1}s)
          +b_{n+1}(a_0s)-b_n(a_1s)
\]
computes $d_nd$: the first two terms use $A_8$, and the last two
use $B_7$. All generator singular products have therefore been
expressed by the same polynomials in the two algebras.

For the uniqueness step, take the vertex subalgebra $D\subset V\times V'$
generated by the diagonal pairs of the four named fields. Common
polynomial OPE closure and filtered Wick induction show that $D$ is
their differential normal span. Every generator has positive weight,
so $D_0=\mathbb C(\mathbf1,\mathbf1)$ and $D_n=0$ for $n<0$.
It follows that $D$ has a unique maximal proper graded ideal: every
proper graded ideal has zero degree-zero part, and so does their sum.
The coordinate projections are surjective because $L,a$ generate
each target. They preserve the conformal vector and hence the grading.
Their kernels are maximal proper graded ideals because $V,V'$ are
simple. The kernels coincide, which gives the asserted isomorphism.
No freeness or independence of normal monomials has been used.
\end{proof}

The exact parafermion records instantiate this proposition at
$c=8/5$ and $\lambda=14/473$. The upper-pole record supplies the
self-row and the remaining positive poles, the nonlinear record
supplies $a_0b$, and the truncation record supplies $a_1d$.
The algebraic simple-pole derivation supplies $a_0d$ from the
independently checked upper poles and skew symmetry. The
Dong--Lam--Yamada change-of-generators record matches the seven
weight-six decoupling coefficients with those of
$K(\mathfrak{sl}_2,13)$, using \cite[Appendix B]{DLY}. The common first rows are the recursive
formulas (5.4), (5.7), and (A.2) of Linshaw \cite{Linshaw}. Simplicity of the actual
generated subalgebra follows from its PCT stability; its translation
operator agrees with the ambient one by the complementary positive
energy argument. Thus the proposition transports the checked
polynomial identities to an identification with the parafermion
algebra. The exact membership, PCT, translation, and polynomial
identities remain the explicit finite inputs to this identification.

\begin{lemma}[Orthogonal energy bound]\label{lem:positive-c2}
Let $V$ be a unitary VOA of CFT type admitting an identity of the form
\eqref{eq:finite-positive-stress}, with the properties just stated.
Then $V$ is $C_2$-cofinite.
\end{lemma}
\begin{proof}
For an ordinary simple $A_i$-module $M$, put
\[
 C_{2,A_i}(M)=\operatorname{span}\{a_{-2}m:a\in A_i,\ m\in M\}.
\]
These quotients are finite-dimensional by \cite[Proposition~5.2]{ABD}.
They are graded for the intrinsic operator $L^{e_i}(0)$.
There are finitely many simple $A_i$-modules. Hence
there is a finite number $D_i$ bounding the intrinsic energy in all
their $C_{2,A_i}$ quotients; this includes their lowest weights.
Regularity decomposes the weak $A_i$-module $V$ into ordinary simples.
Consequently every intrinsic eigenvector of energy above $D_i$ lies
in $C_{2,A_i}(V)$. The intrinsic energy preserves each finite-dimensional
ambient grade $V_n$ and is self-adjoint for its positive form. We may
therefore diagonalize it on $V_n$. Since
\[
 C_{2,A_i}(V)\subset C_2(V),
\]
every vector in $V_n\cap C_2(V)^\perp$ has zero projection to the
intrinsic eigenvectors of energy above $D_i$. It follows that
\[
 \langle L^{e_i}(0)v,v\rangle\le D_i\|v\|^2
 \qquad (v\in V_n\cap C_2(V)^\perp).
\]
For $v\in V_n\cap C_2(V)^\perp$, the conformal-vector identity gives
\[
 n\|v\|^2
 =\sum_i c_i\langle L^{e_i}(0)v,v\rangle
 \le\Bigl(\sum_i c_iD_i\Bigr)\|v\|^2.
\]
The subspace $C_2(V)$ is ambient graded. Its orthogonal complement in
each finite-dimensional $V_n$ represents $(V/C_2(V))_n$, so the last
inequality makes the quotient vanish for
$n>\sum_i c_iD_i$. Finiteness of each $V_n$ proves the claim.
No simultaneous diagonalization of the different energy operators is
required. The same proof applies directly to averages, since every
conjugate factor has the same bound $D_i$.
\end{proof}

\begin{lemma}[Highest spaces for a compatible subalgebra]
\label{lem:compatible-top}
Let $A\subset V$ be grading compatible and regular, and let $M$ be a
weak $V$-module. Suppose a finite-dimensional subspace $E\subset M$ is
stable under every ambient degree-zero mode from $A$ and is killed by
every mode from $A$ of negative ambient degree. Then $E$ is a direct
sum of irreducible Zhu-algebra modules belonging to the tops of simple
ordinary $A$-modules. In particular $L^A(0)|_E$ is diagonalizable,
with eigenvalues among their lowest weights.
\end{lemma}
\begin{proof}
Write
\[
 \Omega_A(M)=\{v\in M:a_nv=0\text{ for all homogeneous }a\in A,
                   \ n\ge\operatorname{wt}_A(a)\}.
\]
Compatibility puts $E$ inside this space. Regularity decomposes
$M\downarrow_A$ as an algebraic direct sum of ordinary simples, and
$\Omega_A$ commutes with that sum. The highest space of an ordinary
simple is its top. Zero modes endow $\Omega_A(M)$ with the corresponding
semisimple Zhu-algebra action; its invariant finite-dimensional
subspace $E$ is therefore a sum of simple tops. This also explains why
one must impose the lowering condition for all states of $A$, rather
than only for a chosen generating set.
\end{proof}

\begin{lemma}[Vacuum-like vectors]\label{lem:vacuum-like}
Let $M$ be a weak module over a VOA $V$. If $v\in M$ satisfies
$L_M(-1)v=0$, then $a_nv=0$ for all $a\in V$ and $n\ge0$, and
$a\mapsto a_{-1}v$ is a $V$-module homomorphism $V\to M$ sending
$\mathbf1$ to $v$.
\end{lemma}
\begin{proof}
Translation covariance gives
$L_M(-1)a_{n+1}v=-(n+1)a_nv$. Lower truncation and downward induction
give $a_nv=0$ for every $n\ge0$. The module Jacobi identity then gives
$b_m(a_{-1}v)=(b_ma)_{-1}v$; equivalently this follows by taking the
constant coefficient in weak associativity, since the fields applied
to $v$ have no negative powers. This is the asserted homomorphism.
\end{proof}

\begin{lemma}[Lowest weights and projectivity of the vacuum]
\label{lem:vacuum-projective}
Let $V$ satisfy Lemma~\ref{lem:positive-c2}, and suppose every component
algebra has nonnegative simple lowest weights and its vacuum is the
unique simple of weight zero. Every simple ordinary $V$-module has
nonnegative lowest conformal weight, the only weight-zero simple is
$V$, and $V$ is projective among grading-restricted generalized modules.
\end{lemma}
\begin{proof}
Let $E=M_h$ be the lowest space of a simple ordinary module.
For a homogeneous state $a$,
$[L_V(0),a_n]=(\operatorname{wt}_V(a)-n-1)a_n$ on every module.
Thus degree-zero modes preserve $E$, and lowering modes kill $E$.
Lemma~\ref{lem:compatible-top} implies that every
$L^{e_i}(0)|_E$ is diagonalizable with nonnegative real spectrum.
Taking traces, rather than assuming the operators commute, gives
\begin{equation}
 h\dim E=\sum_i c_i\operatorname{Tr}_E L^{e_i}(0)\ge0.
 \label{eq:top-trace}
\end{equation}
In particular $h$ is real and nonnegative. If $h=0$, every summand vanishes.
Hence every partial energy is zero
on $E$, which lies in sums of vacuum tops for every $A_i$. Consequently
$L^{e_i}(-1)E=0$, and \eqref{eq:finite-positive-stress} gives
$L_V(-1)E=0$. By Lemma~\ref{lem:vacuum-like}, any nonzero $v\in E$
gives a nonzero module map $V\to M$, hence an
isomorphism.

Now let $\pi:N\twoheadrightarrow V$ be an epimorphism in the generalized
module category. $C_2$-cofiniteness gives finite length, and the preceding
argument excludes negative weights in all its composition factors.
The finite generalized weight-zero space $N_{[0]}$ therefore has no
lower-weight target for a lowering mode. It is again in every
$\Omega_{A_i}(N)$ and is invariant under degree-zero modes. On this
space $L_V(0)$ is nilpotent, so
\[
 0=\operatorname{Tr}_{N_{[0]}}L_V(0)
  =\sum_i c_i\operatorname{Tr}_{N_{[0]}}L^{e_i}(0).
\]
Diagonalizability and nonnegativity of each partial energy imply it is
zero. The vacuum-top argument gives $L_V(-1)N_{[0]}=0$. Generalized
weight decompositions are preserved by $\pi$, so one can choose
$v\in N_{[0]}$ with $\pi(v)=\mathbf1$. Then
\[
 s:V\longrightarrow N,\qquad s(a)=a_{-1}v,
 \qquad \pi s=\mathrm{id}_V,
\]
is a section. Every epimorphism onto the vacuum splits.
\end{proof}

\begin{lemma}[A projective unit in a finite braided $r$-category]
\label{lem:projective-unit}
Let $\mathcal C$ be a finite complex linear abelian braided monoidal
category with simple unit, right-exact tensor product, and
Grothendieck--Verdier duality with dualizing object $\mathbf1$.
If $\mathbf1$ is projective, then $\mathcal C$ is rigid and semisimple.
\end{lemma}
\begin{proof}
Realize $\mathcal C$ as
finite-dimensional modules for a finite-dimensional algebra, denote
underlying vector-space dimension by $d_0$, and let $P$ correspond to
its regular module. Every $Z$ is a quotient of $P^{\oplus d_0(Z)}$.
For fixed $X$, right exactness gives
\[
 d_0(Z\otimes X)\le C_Xd_0(Z),\qquad
 \dim\operatorname{End}(X^{\otimes n})
 \le d_0(\mathbf1)^2C_X^{2n}<n!
\]
for large $n$, where $C_X=\max\{1,d_0(P\otimes X)\}$.
This uses only right exactness. For nonzero indecomposable $X$,
Grothendieck--Verdier evaluation
$DX\otimes X\to\mathbf1$ is nonzero, hence surjective onto the simple
unit, and splits by projectivity. A finite abelian category is
idempotent complete, and its Hom spaces are finite-dimensional.
Thus $X$ is non-negligible and has moderate growth in the precise
sense of \cite[Theorem~1.1]{EP}; that theorem proves $X$ rigid.
Finite direct sums give
rigidity of all objects. Tensoring is now exact, and
\[
 \operatorname{Hom}(X,-)
 \cong\operatorname{Hom}(\mathbf1,-\otimes X^*)
\]
is exact. Every object is projective, so $\mathcal C$ is semisimple.
\end{proof}

\begin{theorem}[Positive stresses imply strong rationality]
\label{thm:positive-stress-rationality}
A simple unitary VOA of CFT type satisfying the hypotheses of
Lemma~\ref{lem:vacuum-projective} is strongly rational. In particular,
the actual generated algebra $U$ is strongly rational, independently
of its properness in $B$.
\end{theorem}
\begin{proof}
Unitarity gives self-contragredience. Lemmas~\ref{lem:positive-c2} and
\ref{lem:vacuum-projective} give $C_2$-cofiniteness and a projective
vacuum. The category of grading-restricted generalized modules is
finite abelian and braided, its tensor product is right exact, and
contragredient duality is Grothendieck--Verdier duality with dualizing
object $V'=V$; see \cite[Sections~2.3 and~3]{McRae}.
Lemma~\ref{lem:projective-unit} proves semisimplicity, hence rationality
in the $C_2$-cofinite setting. This argument does not assume that the
nonvacuum modules are unitarizable.
The local argument is recorded in
\artifact{hi_cft_20260923_next_structure_rationality.md}.
\end{proof}

\section{Branching and the boundary fusion ring}
\label{sec:branching}

From this point we use the properness conclusion and $\dim U_1=6$ of
the finite closure argument. Before that conclusion is supplied, every
statement in this section is conditional on the proper branch.

\subsection{A universal top bound and actual lattice branching}

The sharper bound uses an actual regular subalgebra $A_L\subset U$.
In the notation of the local construction, a rational lattice--Potts
core $\mathcal T$ of central charge $14/5$ and $K(\mathfrak{sl}_2,13)$
of central charge $8/5$ have pointed charges
$(j,k)\in(\mathbb Z/13)^2$ with quadratic form
$-(j^2+k^2)/13$. Put $A=\mathcal T\otimes K(\mathfrak{sl}_2,13)$.
The field certificate supplies nonzero weight-two vectors $v$ with
\[
 (\omega_U-\omega_A)_1v=0,\qquad
 z_1v=\frac{12}{13}v,\qquad
 (\omega_{\mathcal T})_1v=\frac{14}{13}v,
 \qquad p_1v=\frac23v,
\]
at Cartan weights $(\pm1,\pm5)$ and $(0,\pm8)$.
These are the top characters of the simple-current charges
$(5,1)$ and its inverse, up to reversing the second charge.
The label is forced by actual, unreduced lowest weights: the
$\tau$-labelled $\mathcal T$ module of charge $5$ starts at $31/65$,
whereas $14/13-31/65$ is not integral, and the only parafermion modules
starting at $12/13$ are the two inverse simple currents.

Here is why the witness gives an embedded extension. The kernel of
$(\omega_U-\omega_A)_1$ in $U$ is a connected unitary extension of
$A$ with conformal vector $\omega_A$. Regularity of $A$ and finite
ambient homogeneous spaces make this extension finite. In a connected
unitary extension, invertible constituents have multiplicity at most
one and their support is a subgroup: the invariant pairing gives a
nonzero product of a constituent with its inverse into the vacuum
line; associativity then makes multiplication by it invertible on
each supported invertible line. Integer grading makes that support
isotropic. The only isotropic lines for the displayed quadratic form
are $k=5j$ and $k=-5j$, since $5^2=-1\pmod{13}$. The witnessed
nonzero charge generates one of these lines, and both cannot occur
because together they generate nonintegral-spin charges. Consequently,
after the harmless choice of parafermion charge orientation,
\[
 A_L=\bigoplus_{k=-5j}\mathcal T_{1,j}\otimes\mathrm{PF}_{0,k}
 \subset U,\qquad
 \operatorname{Rep}(A_L)\simeq
 \mathrm{Fib}\boxtimes\mathrm{SO}(3)_{13}.
\]
The field witness and extension argument are in
\artifact{hi_cft_20260923_next_structure_order13_core.md} and
\artifact{hi_cft_20260923_next_structure_core_extension.json}.
There are fourteen simple modules, labelled $x_a=(1,a)$ and
$y_a=(\tau,a)$ for $0\le a\le6$. Indeed the free order-thirteen
extension removes the hyperbolic pointed factor; no orbit has a
stabilizer, so there is no fixed-point resolution. The two remaining
labels are the Fibonacci label and the even affine label $2a$.

We give enough charge data to reconstruct their tops. For
$\epsilon\in\{1,\tau\}$ put
\[
 p_1(m)=\begin{cases}0,&3\mid m,\\2/3,&3\nmid m,\end{cases}
 \qquad
 p_\tau(m)=\begin{cases}2/5,&3\mid m,\\1/15,&3\nmid m.\end{cases}
\]
The actual lowest weights of the lattice--Potts modules are
\begin{equation}
 h_{\mathcal T}(\epsilon,j)
 =\min_{m\equiv j\ (13)}
 \left(\frac{m^2}{156}+\frac{\mathbf1_{m\text{ odd}}}{4}
                   +p_\epsilon(m)\right).
 \label{eq:core-charge-minimum}
\end{equation}
An even minimizing $m$ contributes a left affine singlet and an odd
one a left doublet; their right Cartan charge is $m$. Choosing
$|m|\le6$ bounds the minimum by $14/13$, while $|m|\ge13$ gives
$m^2/156>14/13$. Thus $-12\le m\le12$ contains every minimizer.
For the representative $-6\le k\le6$, the parafermion formula is
\begin{equation}
 h_{\mathrm{PF}}(a,k)
 =\frac{a(a+1)}{15}-\frac{k^2}{13}+\max(0,|k|-a).
 \label{eq:pf-charge-minimum}
\end{equation}
The last term is the least affine grade needed for charge $2k$
outside the finite $\mathfrak{sl}_2$ top. This is equivalently the
actual-top formula in \cite[Section~2.4]{ALY}; the parafermion top is
one-dimensional. Finally minimize
$h_{\mathcal T}(\epsilon,j)+h_{\mathrm{PF}}(a,-5j)$ over
$j\in\mathbb Z/13$, using the representative in
\eqref{eq:pf-charge-minimum}.

Let $E$ be any simple $U$-module top and write $x_a,y_a$ also for the
nonnegative integer multiplicities of these $A_L$ tops in $E$.
Lemma~\ref{lem:compatible-top} justifies taking complete $A_L$ tops
here: the finite space $E$ is invariant under all their Zhu zero modes,
not merely contained in the sum of their highest spaces.
Set $D=\dim E$. For each top, let $L$ be the trace of the left affine
energy, $Q=\operatorname{Tr}H_R(0)^2$, and $P,Z$ the Potts and
parafermion energy traces. Its contribution to \eqref{eq:top-trace} is
\begin{equation}
 W=L+\frac{83}{6084}Q+\frac{25}{6}P+\frac{25}{26}Z.
 \label{eq:top-W}
\end{equation}
This is a trace formula on a complete right $\mathfrak{sl}_2$
representation, not an assertion about a single partial-energy
eigenvector. In fact, for an irreducible right $\mathfrak{sl}_2$
representation of highest weight $r$,
\[
 \frac{\operatorname{Tr}H_R(0)^2}{\dim}=\frac{r(r+2)}3,
 \qquad L^{A_{1,39}}(0)=\frac{r(r+2)}{164}.
\]
Thus the Sugawara trace is $3Q/164$, and multiplying by $3403/4563$
gives $83Q/6084$. Haar conjugation preserves traces on $E$.
The exact dimensions and trace contributions are
\[
\begin{array}{c|r|r@{\qquad}c|r|r}
 \text{top}&d&4680W&\text{top}&d&4680W\\\hline
 x_0&1&0&y_0&1&7800\\
 x_1&1&600&y_1&7&36960\\
 x_2&5&10440&y_2&15&110880\\
 x_3&9&56640&y_3&2&4080\\
 x_4&2&5520&y_4&14&81480\\
 x_5&14&148560&y_5&6&18240\\
 x_6&6&54240&y_6&2&4920
\end{array}
\]
For clarity, the lowest core weights and all minimizing absolute
charges used in this table are
\[
\begin{array}{c|c|l@{\qquad}c|c|l}
 &h_{A_L}&|j|&&h_{A_L}&|j|\\\hline
 x_0&0&0&y_0&2/5&0\\
 x_1&2/15&0&y_1&8/15&0,5\\
 x_2&2/5&0,3&y_2&4/5&0,2,3,5\\
 x_3&4/5&0,2,3,6&y_3&1/5&2\\
 x_4&1/3&6&y_4&11/15&1,2,4,6\\
 x_5&1&1,4,6&y_5&2/5&1,4\\
 x_6&4/5&4&y_6&1/5&4
\end{array}
\]
For each positive $j$ both signs occur. A convenient independent
check on every row is
\[
 Z=h_{A_L}d-L-Q/156-P,\qquad
 W=\frac L{26}+\frac{7Q}{936}+\frac{125P}{39}
                              +\frac{25}{26}h_{A_L}d.
\]
Complete right $\mathfrak{sl}_2$ characters imply
$y_2\ge x_5$ and $y_4\ge x_5+x_6$: these compare charges $10,12$
and charges $7,9$, respectively, with the common left doublet removed
in the second comparison. Direct comparison of the displayed table gives
\begin{align}
 4680\bigl(\tfrac{17}{10}D-W_{\rm total}\bigr)
 ={}&8460(y_2-x_5)+28716(y_4-x_5-x_6)\notag\\
 &+7956x_0+7356x_1+29340x_2+14964x_3+10392x_4\notag\\
 &+22212x_6+156y_0+18732y_1+11832y_3\notag\\
 &+1188y_4+29496y_5+10992y_6\ge0.
 \label{eq:sharp-top-bound}
\end{align}
Since $W_{\rm total}=hD$, this proves $0\le h\le17/10<2$ for
\emph{every} simple $U$-module, without assuming it is unitarizable.

\begin{reviewnote}{R2: scope of the core and finite-top certificates}
The charge minimizations and the nonnegative identity
\eqref{eq:sharp-top-bound} are specified above; their finite exact
check is
\artifact{hi_cft_20260923_next_structure_order13_top_check.json}.
The embedded lattice--Potts core and the spectral field witness remain
construction inputs, supplied by the exact core-extension certificate.
The finite-top irreducibility below is likewise imported from exact
field actions, particularly
\artifact{hi_cft_20260923_next_81top_all.json}; it should not be replaced
by numerical cyclicity alone.
\end{reviewnote}

The strongly rational restriction $B\downarrow_U$ is semisimple. Its
simple summands have integer lowest weights, hence only weights $0,1$
are possible. The vacuum occurs once, while the complementary current
space has dimension $86-6=80$. Its certified irreducible zero-mode
action therefore implies
\begin{equation}
 B\downarrow_U=U\oplus X,\qquad h_X=1,\qquad\dim X_{\rm top}=80.
 \label{eq:B-branching}
\end{equation}
The four mixed nontrivial $E_6\oplus A_2$ lattice sectors start at
weight $1$, with $81$-dimensional tops. Their next possible top weight
would be $2$, excluded by \eqref{eq:sharp-top-bound}. Their exact
top-action certificates use a one-dimensional Cartan eigenline,
positive-metric adjoint identities, and a cyclic span of dimension
$81$. The adjoint identities make the generated operator algebra a
star algebra; orthogonal complements are therefore invariant, and
the unique eigenline plus cyclicity proves irreducibility.

Write the actual lattice extension as $H=B\oplus J\oplus J^2$. Its
restriction is consequently
\begin{equation}
 H\downarrow_U=\mathbf1\oplus X\oplus Y\oplus Y^*,\qquad
 h_Y=1,\quad\dim Y_{\rm top}=81.
 \label{eq:H-branching-fork}
\end{equation}
At this stage $Y\cong Y^*$ has not been assumed. The three types
$\mathbf1,X,Y$ are distinguished by their lowest spaces.

\subsection{The rank-four or rank-six alternative}

Put $\mathcal C=\operatorname{Rep}(U)$ and
$\mathcal F=\mathcal C_H$, the category of \emph{all} right modules
over the extension algebra $H$. The strict positivity of nonvacuum
lowest weights gives the positive spherical structure through character
asymptotics. Since $H$ is holomorphic, it is Lagrangian and the
center correspondence \cite[Section~4]{DMNO}
gives $\mathcal C\simeq Z(\mathcal F)$.
Write $s=\operatorname{FPdim}_{\mathcal C}B$. Extension dimension
identities give
\[
 d_X=s-1,\qquad d_Y=d_{Y^*}=s,\qquad
 \operatorname{Dim}\mathcal C=9s^2,\quad
 \operatorname{Dim}\mathcal F=3s.
\]
Pseudo-unitarity of $\mathcal F$ follows from that of its center and
the squared global-dimension and Frobenius--Perron-dimension identities.
No $C^*$ structure is needed in this section. All categorical dimensions
below use the positive spherical structure and therefore equal
Frobenius--Perron dimensions.

For center induction $I$ and forgetful functor $F$, one has $H=I(\mathbf1)$
and $FI(\mathbf1)=\bigoplus_{a\in\operatorname{Irr}\mathcal F}a\otimes a^*$.
Adjunction therefore gives
\begin{equation}
 \operatorname{rank}\mathcal F
 =\dim\operatorname{End}_{\mathcal C}(H)
 =\begin{cases}4,&Y\not\cong Y^*,\\6,&Y\cong Y^*.
 \end{cases}
 \label{eq:rank-fork}
\end{equation}
\begin{lemma}[Automorphisms give invertible boundary objects]
\label{lem:automorphisms-invertibles}
For a pseudo-unitary fusion category $\mathcal F$, algebra
automorphisms of its canonical Lagrangian algebra $I(\mathbf1)$ inject
into $\operatorname{Inv}(\mathcal F)$.
\end{lemma}
\begin{proof}
Normalize the special Frobenius structure so that
$m\Delta=\mathrm{id}$ and set $f*g=m(f\otimes g)\Delta$.
If $f$ is an algebra automorphism, then $f*f=f$, and the same holds
for its composition inverse. Under the two-product description
\cite[Lemma~3.1 and Corollary~3.4]{BJ}, composition is fusion-ring
multiplication and the primitive convolution idempotents are
$d_a[a]$. Thus both $f$ and $f^{-1}$ have expansions with coefficients
$0$ or $d_a>0$. Their fusion product is $[\mathbf1]$.
Nonnegative fusion coefficients rule out cancellation. Every pair of
objects in the two supports must therefore multiply to copies of the
unit; rigidity then forces each pair to consist of inverse invertible
objects. The two supports are singletons and their dimensions are
one. Composition preserves this assignment, proving the injection.
\end{proof}

The actual order-three automorphisms of $H/B$ consequently inject into
$\operatorname{Inv}(\mathcal F)$. This injection is already spherical:
the proof used positivity of dimensions and fusion coefficients,
without using adjoints of morphisms.

\begin{lemma}[No sign extension]\label{lem:no-sign-extension}
The nonvacuum summand $X$ of $B\downarrow_U=U\oplus X$ is not
invertible.
\end{lemma}
\begin{proof}
If $X$ were invertible, self-duality of $B$ and connectedness would
give $X^2=\mathbf1$. The extension would have its $\mathbb Z/2$
grading, and its sign automorphism $g$ would satisfy $B^g=U$.
On currents this would give
$(E_6\oplus A_2)^g=U_1\cong\mathfrak{sl}_2\oplus\mathfrak{sl}_2$,
of rank two.

Every finite-order automorphism of a complex simple Lie algebra is
conjugate to a toral automorphism times a diagram automorphism;
see \cite[Section~5.2]{RLYG}. It therefore fixes a toral subalgebra
whose dimension is the number of orbits of the diagram automorphism
on simple roots. In type $E_6$ this number is at least four, and in
type $A_2$ at least one. Since the nonisomorphic ideals $E_6,A_2$
are preserved separately, $(E_6\oplus A_2)^g$ contains a toral
subalgebra of dimension at least five. Such a subalgebra cannot lie
in the rank-two Lie algebra $U_1$, a contradiction.
\end{proof}

If rank six had six invertibles, then $3s=6$ and
$d_X=s-1=1$, contradicting Lemma~\ref{lem:no-sign-extension}.
Thus there are
exactly three invertibles in either rank. Their associator is also
determined: the local $B$-module category has metric group
\[
 (\mathbb Z/3)^2,\qquad q(a,b)=\frac{2a^2+b^2}{3}\pmod1.
\]
The actual extension $H/B$ selects the diagonal isotropic line
$L_+=\{(a,a)\}$. Its complement
$L_-=\{(a,-a)\}$ is also isotropic, and
$(\mathbb Z/3)^2=L_+\oplus L_-$.
The pointed subcategory on $L_-$ is symmetric with trivial twists,
hence Tannakian; as a tensor category it is the untwisted
$\operatorname{Vec}_{\mathbb Z/3}$. Free $H/B$-module induction is
tensor on this subcategory. It is fully faithful since
\[
 \operatorname{Hom}_{H/B}(x\otimes(H/B),y\otimes(H/B))
 \cong\operatorname{Hom}_{\operatorname{Rep}(B)}
       (x,y\otimes(H/B))
\]
has dimension $\delta_{x,y}$ for $x,y\in L_-$: distinct elements of
$L_-$ represent distinct cosets of $L_+$. Every free module is a
direct sum of these three simple modules, and separability makes
every module a summand of a free module. Thus induction identifies the entire
all-module category for $H/B$ with this untwisted pointed category.
Viewing these modules in $\operatorname{Rep}(U)$ forgets only the
locality condition for the intermediate algebra $B$; it preserves
$H$-module morphisms and relative tensor products. This gives the
required full tensor embedding into $\mathcal F$.
Merely knowing the group of invertibles would not determine this
associator.

\subsection{Excluding rank four}

At rank four the fusion ring is near-group:
\[
 a^3=1,\qquad a\rho=\rho a=\rho,\qquad
 \rho^2=1+a+a^2+m\rho.
\]
There are exactly two self-dual simples, $\mathbf1$ and $\rho$.
Consequently \cite[Theorem~1.1]{Larson}, which assumes
pseudo-unitarity, restricts $m$ to $0,2,3,6$.
For $m=0$, $s=2$ contradicts Lemma~\ref{lem:no-sign-extension}.
For $m=2$, the categorical reduction
\cite[Theorem~7.1 and Corollary~7.4]{EGO} supplies the twisted $A_4$ double.
Its electric regular algebra gives a holomorphic $c=8$ extension
$H'$ with $U=(H')^{A_4}$ by the extension correspondence
\cite{HKL}. The nonvacuum lowest weights are positive, so $H'$ is
of CFT type, and \cite[Theorem~1]{DM} identifies it with $V_{E_8}$.
For completeness, the nontrivial order-two fixed Lie algebras are
$D_8$ and $E_7+A_1$, with dimensions $120,136$ and adjoint traces
$-8,24$. At order three the possibilities are
$A_8$, $E_6+A_2$, $D_7+\mathbb C$, and $E_7+\mathbb C$,
with dimensions $80,86,92,134$ and traces $-4,5,14,77$.
These lists follow by selecting Kac coordinates of total mark two
or three in the affine $E_8$ diagram, whose marks are
$1,2,2,3,3,4,4,5,6$. For order three the trace is
$(3\dim\mathfrak g^g-248)/2$, since the two nontrivial eigenspaces
are dual. Thus the two relevant traces are bounded below by $-8$
and $-4$. Averaging over $A_4$ gives
\[
 \dim(E_8^{A_4})\ge\frac{248+3(-8)+8(-4)}{12}=16,
\]
contradicting $\dim U_1=6$.

The remaining cases require a spherical induction argument before any
comparison with affine characters. We give this reduction explicitly.

\begin{lemma}[Near-group induction]\label{lem:near-group-induction}
Suppose $m=3k$, and put $d^2=3+3kd$, $s=2+kd$.
There are integers $p,q\ge0$, $p+q=k$, and a primitive cube root
$\zeta$ such that the known center objects have forgetful images
\[
\begin{array}{c|c|c}
 \text{object}&F(\text{object})&\text{twist}\\\hline
 X&1+k\rho&1\\
 Y,Y'&1+a+k\rho,\ 1+a^2+k\rho&1,1\\
 C&a+a^2+k\rho&\zeta\\
 D,D'&a+p\rho,\ a^2+p\rho&\zeta^2,\zeta^2\\
 E,E'&a+q\rho,\ a^2+q\rho&\zeta^2,\zeta^2
\end{array}
\]
Every further simple $Z_j$ forgets to $n_j\rho$, $n_j>0$. Its twist
$\theta_j$ satisfies
\begin{align}
 \sum_jn_j^2&=6+3k^2+4pq,\label{eq:near-group-length}\\
 S_1:=\sum_jn_j^2\theta_j&=k^2-kd-(p-q)^2\zeta^2,\notag\\
 \sum_jn_j^2\theta_j^2&=\overline{S_1}
                 +\varepsilon\sqrt{9k^2+12},
 \qquad\varepsilon=\nu_2(\rho)\in\{1,-1\}.
 \label{eq:near-group-moments}
\end{align}
\end{lemma}
\begin{proof}
The four one-dimensional representations of the fusion algebra have
formal codegrees $3s,9s/d^2,3,3$. Thus induction of the unit is
$1+X+Y+Y'$, with dimensions $1,d^2/3,s,s$ and trivial twists.
Since $FI(1)=4\cdot1+a+a^2+3k\rho$, adjunction and comparison of
rational and irrational dimension coefficients give the first two
rows. Also
\[
 FI(a)=1+4a+a^2+3k\rho,\qquad
 FI(a^2)=1+a+4a^2+3k\rho.
\]
The two inductions have norm four, share one constituent, and contain
$Y,Y'$ respectively. Their remaining norm-three parts are therefore
multiplicity-free, with a single common constituent $C$.
Write its $\rho$ multiplicity as $c$ and the other two as $p,q$.
Then $c+p+q=2k$.

The untwisted pointed subcategory gives $\nu_3(a)=1$, so positivity
in the twist-trace formula forces every constituent of $I(a)$ to
have cube-root twist. Its first indicator, separated over the
disjoint fields $\mathbb Q(d)$ and $\mathbb Q(\zeta)$, gives
\[
 2+2\theta_C+\theta_D+\theta_E=0,\qquad
 k+c\theta_C+p\theta_D+q\theta_E=0.
\]
The first equality forces $\theta_C=\zeta$ and
$\theta_D=\theta_E=\zeta^2$, up to simultaneous conjugation; the
second gives $c=k$ and $p+q=k$. Ribbon duality preserves the primed
twists. Subtracting the nine known squared dimensions from $9s^2$
gives \eqref{eq:near-group-length}. The first two twist traces of
$I(\rho)$ give \eqref{eq:near-group-moments}, using
$3/d=d-3k$ and $3s/d=\sqrt{9k^2+12}$.
\end{proof}

\begin{proposition}[The multiplicity-three modular representation]
\label{prop:m3-modular}
For $k=1$, the center has the modular data of
$G_{2,3}\otimes A_{2,1}$ or their simultaneous reversal.
\end{proposition}
\begin{proof}
One of $p,q$ is zero. The corresponding $D,D'$ are invertible and
generate a nondegenerate pointed category $\mathcal P$ with twists
$1,\zeta^2,\zeta^2$. Write
$Z(\mathcal F)=\mathcal M\boxtimes\mathcal P$ using its modular
centralizer. The known objects form three $\mathcal P$-orbits;
representatives $1,X,C$ in $\mathcal M$ have dimensions
$1,1+d,s$ and twists $1,1,\zeta$. Balancing verifies their trivial
monodromy with $D$ from $D X=E$ and $DC\in\{Y,Y'\}$.
The remaining free orbits, together with
\eqref{eq:near-group-length}, leave three simples in $\mathcal M$,
all of dimension $d$.

Take $\zeta=e^{2\pi i/3}$, the other choice conjugating the answer.
For their twists $z_1,z_2,z_3$, the two moments become
\[
 \sum_i z_i=\alpha:=\frac{-1-i\sqrt7}{2},\qquad
 \sum_i z_i^2=-\frac12-i(\varepsilon-\tfrac12)\sqrt7.
\]
The second sum would have absolute value four for $\varepsilon=-1$,
so $\varepsilon=1$. Since $\alpha^2+\alpha+2=0$, Newton's identity
and the unit modulus of all $z_i$ give elementary symmetric functions
$e_2=\overline\alpha$, $e_3=1$. Thus the twists are the roots of
$t^3-\alpha t^2+\overline\alpha t-1$, namely
$\zeta_7^3,\zeta_7^5,\zeta_7^6$.

The actual extension algebra $1+X$ lies in $\mathcal M$. Its local
module category has dimension three and is the pointed category with
nontrivial twist $\zeta$, fixed by its multiplicative central charge.
Its two nonvacuum simples must both restrict to $C$: each underlying
object has dimension $s$, cannot contain the unit by induction
reciprocity, and no other combination of the six known dimensions
has these properties. Hence the branching columns are $1+X,C,C$.
Modular intertwining for this extension
\cite[Theorems~4.1 and~4.5]{KO} fixes the first three rows
of $\widehat S=\sqrt3sS_{\mathcal M}$:
\[
 \begin{pmatrix}
 1&1+d&s&d&d&d\\
 1+d&1&s&-d&-d&-d\\
 s&s&-s&0&0&0
 \end{pmatrix}.
\]
All six objects are self-dual, the last three by their distinct
twists. Their fusion matrices are real symmetric, so Verlinde
diagonalization and the positive vacuum row make $S_{\mathcal M}$
real. For the remaining block $K$, orthogonality and the modular
relation yield
\[
 \sum_jK_{ij}=d,\qquad
 \sum_jK_{ij}z_j=d+i\sqrt3s\,z_i^{-1}.
\]
These determine each real row: three distinct points on a circle are
not collinear. The affine $G_{2,3}$ data
\cite[Table~4 and equation~(3.8)]{AGJS} realize these rows and twists
in the reversed orientation; conjugating gives the other possibility.
Thus the entire representation, rather than merely its twist list,
is forced.
\end{proof}

\begin{lemma}[The balanced multiplicity-six pattern is impossible]
\label{lem:m6-balanced}
In Lemma~\ref{lem:near-group-induction} with $k=2$, one has
$\{p,q\}=\{0,2\}$.
\end{lemma}
\begin{proof}
Suppose $p=q=1$, and put $d=3+2\sqrt3$, $r=1+d$.
The known dimensions are $1$, $1+2d$, three copies of $2r$, and four
copies of $r$. The remaining dimensions are $n_jd$, with
$\sum_jn_j^2=22$. Write
$u=(2-\sqrt3)/12$, $v=(2+\sqrt3)/12$.
The normalized vacuum row is
\[
 (u,v,1/3,1/3,1/3,1/6,1/6,1/6,1/6,n_j\sqrt3/12).
\]
Signed Galois symmetry \cite[Theorem~II(iii)]{DLN}, for an
automorphism reversing $\sqrt3$, sends this row to the $X$ row:
its first entry becomes $v>0$, uniquely identifying $X$ and a
positive sign. The forced row is
\[
 (v,u,1/3,1/3,1/3,1/6,1/6,1/6,1/6,-n_j\sqrt3/12).
\]
The object $X$ is self-dual because its dimension is unique.
Verlinde's formula consequently gives
\[
 N_{XX}^{X}=\frac{v^3}{u}+\frac{u^3}{v}
              +\frac39+\frac4{36}-\frac{22}{48}
            =\frac{97}{72}+\frac{32}{72}-\frac{33}{72}
            =\frac43,
\]
contradicting integrality. Here $u+v=1/3$, $uv=1/144$ verify the
first two terms directly.
\end{proof}

\begin{lemma}[A universal multiplicity-six character projection]
\label{lem:m6-projection}
The unbalanced case has a nondegenerate pointed factor $\mathcal P$
with twists $1,\zeta^2,\zeta^2$. In its modular centralizer there is
a character projection with matrices
\begin{equation}
 S_v=-\frac12\begin{pmatrix}1&2&1\\1&0&-1\\1&-2&1\end{pmatrix},
 \qquad T_v=\operatorname{diag}(1,-i,-1),
 \label{eq:m6-projection}
\end{equation}
before the common central-charge scalar, or the conjugate matrices.
\end{lemma}
\begin{proof}
The zero multiplicity gives the pointed factor just as for $k=1$.
In the centralizer the known simples have dimensions
$1,2d+1,2d+2$ and twists $1,1,\zeta$; all others have dimensions
$n_jd$, with $\sum n_j^2=6$.
The Gauss sum and the second indicator give
\[
 \sum n_j^2z_j=-2-4i,\qquad
 \sum n_j^2z_j^2=-2+4i(1-\varepsilon).
\]
The second equality excludes $\varepsilon=-1$ by absolute values.
Since $(1+\Re z)(1+\Im z)\ge0$ on the unit circle, summing gives
$6-2-4+\tfrac12\Im(-2)=0$. Thus each $z_j$ is $-1$ or $-i$,
with total squared multiplicity two or four, respectively.

Put
\[
 v_0=\chi_1-\chi_X,\qquad
 v_A=-\frac12\sum_{z_j=-i}n_j\chi_j,\qquad
 v_B=-\sum_{z_j=-1}n_j\chi_j.
\]
The vacuum row is $(u,v,1/\sqrt3,n_j/4)$, with
$u=(2\sqrt3-3)/12$, $v=(2\sqrt3+3)/12$; Galois symmetry gives
the $X$ row $(v,u,1/\sqrt3,-n_j/4)$.
Subtracting these rows yields
$Sv_0=-v_0/2-v_A-v_B/2$. All three combinations are fixed by
duality. Applying $S^2=1$ and
$STS=iT^{-1}ST^{-1}$ to $v_0$ then determines
\[
 Sv_A=-v_0/2+v_B/2,\qquad
 Sv_B=-v_0/2+v_A-v_B/2,
\]
which is \eqref{eq:m6-projection}. This proof permits either one
dimension-$2d$ object or four dimension-$d$ objects at twist $-i$;
no unproved rank-nine classification has been used.
\end{proof}

\begin{lemma}[Pole cancellation]\label{lem:pole-cancellation}
Let $\sigma$ be a finite-image representation of
$\mathrm{SL}_2(\mathbb Z)$ with no invariant vector. Two weight-zero
vector-valued modular functions, holomorphic on the upper half-plane,
with representation $\sigma$ and equal principal parts at infinity
are equal.
\end{lemma}
\begin{proof}
Their difference has no pole at infinity. Covariance makes every
component holomorphic at every cusp of the finite-index kernel of
$\sigma$. It descends to a holomorphic function on the compactified
modular curve and is constant. The resulting constant vector must be
invariant under $\sigma$, hence zero.
\end{proof}

\begin{proposition}[Current obstructions for the remaining near groups]
\label{prop:near-group-current-obstructions}
Under the hypotheses above, $m=3$ forces
\[
 \dim U_1=15+A+10B\quad\text{or}\quad14+42A,
 \qquad A,B\in\mathbb Z_{\ge0},
\]
whereas $m=6$ forces
\[
 \dim U_1=8+A\quad\text{or}\quad5+12P,
 \qquad A,P\in\mathbb Z_{\ge0}.
\]
In particular none allows $\dim U_1=6$.
\end{proposition}
\begin{proof}
We give the theta normalizations, since a factor of two in the vacuum
coordinate would change these obstructions. For $m=3$, use the affine
$G_{2,3}$ labels $1,7,27,77,14,64$, and for each $A_{2,1}$ label put
\[
 v=(\chi_1-\chi_{64},-\chi_{77},-\chi_{14},-\chi_7).
\]
The complement of the extension character is the even Fourier block
on $\mathbb Z/7$, with entries
\[
 S_{00}=-1/\sqrt7,\quad S_{0r}=-2/\sqrt7,\quad
 S_{r0}=-1/\sqrt7,\quad
 S_{rs}=-2\cos(4\pi rs/7)/\sqrt7,
\]
and twist exponents $r^2/7$, $0\le r\le3$.
Together with the even $A_2$ block, multiplication by $\eta^8$
identifies it with the Weil representation of a positive rank-four
lattice with discriminant form
$\mathbb Z_7(r^2/7)\oplus\mathbb Z_3(s^2/3)$.
The actual extension gives $\chi_1+\chi_{64}=\chi_B$. Thus a
difference of two such character vectors has vacuum projection
$\delta v_0=2\delta\chi_U$.

Take $L=L_7\oplus A_2$, where $L_7$ has Gram matrix
$\left(\begin{smallmatrix}2&1\\1&4\end{smallmatrix}\right)$.
For a basis $u,v$ of $L_7$, write $w=v-u/2$ and $x=pu+tw$.
Set $a=p^2$, $b=7t^2/4$, and $c=(y,y)/2$ on the $A_2$ factor.
The quadratics $H_1=a-b$ and $H_2=a+b-c$ are harmonic.
Poisson summation gives vector-valued theta series of weight four.
Their required coefficients are
\[
\begin{array}{c|c|r|r}
 (r,s)&q\text{-power}&\Theta_{H_1}&\Theta_{H_2}\\\hline
 (1,0)&1/7&-1/7&1/7\\
 (3,0)&2/7&3/7&4/7\\
 (0,0)&1&2&-4
\end{array}
\]
The first row comes from $2w/7$, the second from the two shortest
vectors with $a=1/4,b=1/28$, and the last from two $L_7$ roots and
six $A_2$ roots. The only nonvacuum polar top dimensions are $A$ at
$2/7$ and $B$ at $1/7$, with zero allowed for an integer-shifted top.
Subtract the affine reference, which has $22$ currents and $A=7,B=0$.
Pole matching gives
\[
 \eta^8\delta v=\lambda\Theta_{H_1}+\mu\Theta_{H_2},\qquad
 \lambda=7-A+4B,\quad\mu=7-A-3B.
\]
This identity follows from Lemma~\ref{lem:pole-cancellation}: after
division by $\eta^8$, the remainder has no poles; its only possible
$T$-fixed line is $(r,s)=(0,1)$, which the Fourier $S$ matrix does
not preserve. The vacuum coefficient therefore gives
$2(\dim U_1-22)=2\lambda-4\mu$, the first formula.

For the reverse orientation, realize the negative discriminant form
by the positive lattice with Gram matrix
\[
 G_7^-=
 \begin{pmatrix}2&-1&0&0\\-1&2&-1&0\\
 0&-1&2&-1\\0&0&-1&6\end{pmatrix}.
\]
Its dual generator $\ell=(16,11,6,1)/21$ has energy $8/21$;
$(r,s)\mapsto(9r+7s)\ell$ identifies the discriminant form.
Project the theta series to functions even separately in $r,s$.
Let $H=\|x_{A_3}\|^2-3\|x_{A_3^\perp}\|^2$ and
$F=-6D_2\Theta$, where $D_2=q\,d/dq-E_2/6$ is the modular derivative.
The necessary coefficients after projection are
\[
\begin{array}{c|c|r|r}
 (r,s)&q\text{-power}&F&\Theta_H\\\hline
 (0,0)&0&1&0\\
 (0,0)&1&-84&24\\
 (2,1)&2/21&3/14&-2/7
\end{array}
\]
There are twelve roots, all in $A_3$. The polar coset has shortest
vector $(1,2,3,4)/21$, perpendicular to $A_3$; the separately even
projection halves its coefficient, since the paired coset has no
vector of the same norm. If $A$ is its actual polar top dimension,
matching the coefficient $-A$ gives
$\eta^8v=F+(3+14A)\Theta_H/4$. Here $T$ has no fixed vector, so
pole cancellation again proves equality. Since
$\eta^8\chi_B=1+78q+\cdots$, its vacuum coefficient is
$2\dim U_1-94$. Thus
$2\dim U_1-94=-84+6(3+14A)$, giving $14+42A$.

For $m=6$, tensor the projection in
Lemma~\ref{lem:m6-projection} with its pointed factor and first take
the orientation with twists $1,i,-1$. After multiplication by
$\eta^8$ this is the Weil representation for
$L=A_1\oplus A_1\oplus A_2$, restricted by interchange of the two
$A_1$ factors and conjugation of the $A_2$ charge.
The harmonic quadratic
$H=\|x_{A_1^2}\|^2/2-\|y_{A_2}\|^2/2$ has coefficients
\[
 [q^{1/4}]\Theta_{H,(10,0)}=\frac12,
 \qquad [q]\Theta_{H,(00,0)}=-2.
\]
The first counts the two half-roots in a nontrivial $A_1$ coset;
the second counts four $A_1$ roots and six $A_2$ roots.
Let $A=\sum n_j\dim(M_j)_{\rm top}$ over the quarter-twist polar
modules. This is a nonnegative integer, and its projected polar
coefficient is $-A/2$. Use as reference the order-three simple-current
extension of $A_{2,9}$ tensored with $A_{2,1}$: it has $16$ currents
and $A=8$. Its projected representation agrees with
\eqref{eq:m6-projection}; its matrices are also given in
\cite[Lemma~3.4 and Proposition~3.5]{Yu}. Thus cancellation gives
\[
 \eta^8\delta v=(8-A)\Theta_H,
 \qquad 2(\dim U_1-16)=-2(8-A).
\]
As in the first orientation above, the sole $T$-fixed coordinate is
not $S$-invariant. This gives $\dim U_1=8+A$.

For the reverse orientation use the positive lattice with Gram matrix
\[
 G_6^-=
 \begin{pmatrix}2&-1&-1&-1\\-1&2&0&0\\
 -1&0&2&0\\-1&0&0&4\end{pmatrix}.
\]
Its first three basis vectors span $A_3$; if the basis is $a,b,c,e$,
then $w=e+a+(b+c)/2$ is perpendicular to $A_3$ and has norm three.
The discriminant generators $(b+e)/2,(c+e)/2,2w/3$ have quadratic
values $3/4,3/4,2/3$ and mutually integral pairings. The polar coset
is represented by $-w/3$, with energy $1/6$. This shortest vector
is unique, and the projection does not halve its coefficient.
With $H=\|x_{A_3}\|^2-3\|x_{A_3^\perp}\|^2$ and $F=-6D_2\Theta$,
\[
\begin{array}{c|c|r|r}
 \text{coordinate}&q\text{-power}&F&\Theta_H\\\hline
 \text{vacuum}&0&1&0\\
 \text{vacuum}&1&-84&24\\
 \text{polar}&1/6&0&-1
\end{array}
\]
Let $P$ be the weighted top dimension in the half-twist coordinate
with nonvacuum pointed charge; the weighting is by the integers $n_j$.
Contragredience equates the two pointed-charge sums, so averaging
does not change this integer. Pole cancellation gives
$\eta^8v=F+P\Theta_H$: there is no $T$-fixed coordinate in this
orientation. Comparing the vacuum coefficient yields
$2\dim U_1-94=-84+24P$, hence $\dim U_1=5+12P$.
\end{proof}

The representations used in the last proof have finite image because
they are Weil representations, with the finite-order eta multiplier
removed. Thus the proof does not assume a dimension formula for a
space of vector-valued modular forms. The coefficients $A,B,P$ are
weighted dimensions of actual polar tops, and distinct modules need
not have linearly independent characters.

\begin{remark}[Scope of the rank-four exclusion]
The universal induction, balanced-pattern obstruction, weighted
projection, and theta coefficient normalizations are printed above.
Their source calculations are
\artifact{hi_cft_20260923_next_m3_universal_data.md},
\artifact{hi_cft_20260923_next_structure_m6_universal_projection.md},
\artifact{hi_cft_20260923_next_structure_m3_character_obstruction.md},
and \artifact{hi_cft_20260923_next_structure_m6_character_obstruction.md}.
Lemma~\ref{lem:no-sign-extension} removes the formerly used pointwise
stabilizer calculation $\operatorname{Aut}(B/U)=1$ from this proof.
No classification restricted to $C^*$ categories is used here.
\end{remark}

\subsection{The rank-six ring and exclusion of higher multiplicity}

Rank four being excluded, \eqref{eq:rank-fork} gives
$H=\mathbf1\oplus X\oplus2Y$. The composition algebra of
$\operatorname{End}_{\mathcal C}(H)$, identified with the complex
fusion algebra by \cite[Corollary~3.4]{BJ}, is
\[
 K_0(\mathcal F)\otimes\mathbb C
 \cong\mathbb C\oplus\mathbb C\oplus M_2(\mathbb C),
\]
In particular it is noncommutative. Let $a$ generate the three
invertibles. A permutation of order three on three noninvertibles
is either transitive or trivial. If the left action were trivial,
duality would make the right action trivial as well. The two nontrivial
character idempotents of $\mathbb C[\mathbb Z/3]$ would then be
one-dimensional direct ideals. Their four-dimensional complementary
ideal would have to be $M_2(\mathbb C)$, but the dimension character
is nonzero on that ideal, a contradiction.

Choose a noninvertible $\rho$. The right action commutes with the free
left cyclic action, so it is either the same cycle or its inverse.
The same cycle would make the fusion ring generated by two commuting
elements $a,\rho$, which is impossible. Thus $\rho a=a^2\rho$.
If $\rho^*=a^k\rho$, then
$(a^i\rho)^*=a^{i+k}\rho$; involutivity gives $2k=0$ in
$\mathbb Z/3$, hence $k=0$. All three noninvertibles are self-dual.
Frobenius reciprocity gives one copy of the unit and no other
invertible in $\rho^2$. Write its remaining coefficients as
$m_0,m_1,m_2$. The identity $\rho^2a=a\rho^2$ cyclically equates
these coefficients, so
\begin{equation}
 a^3=1,\qquad\rho a=a^2\rho,\qquad
 \rho^2=1+m(\rho+a\rho+a^2\rho),\qquad m\ge1.
 \label{eq:rank6-ring}
\end{equation}

The following induction calculation determines the multiplicity.
Put
\[
 d=\frac{3m+\sqrt{9m^2+4}}2,\quad s=1+d^2,\quad
 \Delta=9m^2+4,\quad R=\rho+a\rho+a^2\rho.
\]
\begin{lemma}[Center induction and two moments]
\label{lem:rank6-induction}
For a pseudo-unitary categorification of \eqref{eq:rank6-ring} whose
pointed subcategory is untwisted, center induction satisfies
\[
 I(1)=1+X+2Y,\quad d_X=d^2,\quad d_Y=s,
 \qquad I(a)=I(a^2)=C_0+C_1+C_2,
\]
where $\theta_X=\theta_Y=1$, $d_{C_j}=s$, and the $C_j$ twists are
$1,\zeta_3,\zeta_3^2$. Every other central simple $Z_j$ satisfies
$F(Z_j)=n_jR$, with $n_j\in\mathbb Z_{>0}$, and
\begin{align}
 \sum_jn_j^2&=4m^2+2,\notag\\
 \sum_jn_j^2\theta_j&=-\frac{m^2+m\sqrt\Delta}{2},
 \label{eq:indicator-first}\\
 \sum_jn_j^2\theta_j^2
 &=\frac{-m^2+(2\varepsilon-m)\sqrt\Delta}{2},
 \qquad\varepsilon=\nu_2(\rho)\in\{1,-1\}.
 \label{eq:indicator-second}
\end{align}
\end{lemma}
\begin{proof}
The complex fusion algebra has two characters
$a\mapsto1$, $\rho\mapsto d,-d^{-1}$, and the two-dimensional
representation
\[
 a\longmapsto\begin{pmatrix}\zeta_3&0\\0&\zeta_3^2\end{pmatrix},
 \qquad \rho\longmapsto\begin{pmatrix}0&1\\1&0\end{pmatrix}.
\]
Their formal codegrees are respectively $3s,3s/d^2,3$.
The induction--formal-codegree correspondence
\cite[Theorems~2.5 and~2.13]{Ostrik13} gives
$I(1)=1+X+2Y$, with the stated dimensions and trivial twists.
Conjugation by $a$ cycles the three noninvertibles, so their
multiplicities in the forgetful image of a central object coincide.
Since
\[
 FI(1)=6\cdot1+3mR,
\]
adjunction and dimensions give $F(X)=1+mR$ and $F(Y)=2\cdot1+mR$.

Similarly,
\[
 FI(a)=FI(a^2)=3a+3a^2+3mR.
\]
Thus the two induced objects have self-inner-products and mutual
inner-product equal to three. They are equal, multiplicity-free sums
of three simples $C_j$, disjoint from $I(1)$. Adjunction gives
$F(C_j)=a+a^2+b_jR$, where $b_j\ge0$ and $\sum_jb_j=3m$.
Untwistedness gives $\nu_3(a)=1$. The center-twist formula
\cite[Theorem~4.1]{NS} therefore reads
\[
 \sum_jd_{C_j}\theta_{C_j}^3=\operatorname{Dim}\mathcal F
                         =\sum_jd_{C_j}.
\]
Positive dimensions and equality in the triangle inequality imply
$\theta_{C_j}^3=1$. The first indicator of $a$ is zero, so
\[
 \sum_j(2+3b_jd)\theta_{C_j}=0.
\]
The real quadratic field $\mathbb Q(d)$ is disjoint from
$\mathbb Q(\zeta_3)$. Equating its two coefficients shows that the
three twists are $1,\zeta_3,\zeta_3^2$ once each and $b_j=m$ for all
$j$. All invertible constituents have now been exhausted by induction
of $1,a,a^2$; every remaining forgetful image is $n_jR$.

The center dimension identity is
\[
 9s^2=1+d^4+4s^2+9d^2\sum_j n_j^2,
\]
which simplifies to the first asserted equation. Finally adjunction
gives
\[
 I(\rho)=mX+mY+m(C_0+C_1+C_2)+\sum_jn_jZ_j.
\]
Its first and second twist traces are
$\operatorname{Dim}\mathcal F\,\nu_1(\rho)=0$ and
$\operatorname{Dim}\mathcal F\,\nu_2(\rho)=3s\varepsilon$.
Using $d^{-1}=d-3m$ gives \eqref{eq:indicator-first} and
\eqref{eq:indicator-second}.
\end{proof}

The roots in these equations form a multiset of exactly $4m^2+2$
elements, with each $\theta_j$ repeated $n_j^2$ times. No distinctness
of the twists has been assumed.

Write $\Delta=t^2D$ with $D$ squarefree. Since
$(3m)^2<\Delta<(3m+1)^2$, we have $D>1$. Reduction modulo three and
seven gives $D\equiv1\pmod3$ and $7\nmid D$: the latter uses that
$m^2\equiv5\pmod7$ has no solution.
Except when $D=10$,
\begin{equation}
 \frac{\varphi(2D)}{\sqrt D}\ge\frac{12}{\sqrt{13}}.
 \label{eq:squareclass-bound}
\end{equation}
Indeed for odd $D$ the left side is the product of
$(p-1)/\sqrt p$ over its prime factors. A single prime must be at
least $13$, and two distinct primes contribute at least
$40/\sqrt{55}>12/\sqrt{13}$; all further factors exceed one.
For even $D$ there is an additional factor $\sqrt2$.
A single odd prime is $5$ (the exceptional $D=10$) or at least $11$,
and two odd primes again give a larger product.
Apply \cite[Theorem~1.3]{Larson} to two copies of the root multiset in
\eqref{eq:indicator-first}, to remove any half-integral coefficients.
Together with \eqref{eq:squareclass-bound} this gives
\[
 4m^2+2\ge\frac{mt}{2}\varphi(2D)
 \ge\frac{6m\sqrt{9m^2+4}}{\sqrt{13}}.
\]
For $m>1$ this is impossible, since
\[
 36m^2(9m^2+4)-13(4m^2+2)^2
 =4(m^2-1)(29m^2+13)>0.
\]
The exceptional squareclass requires the following equality statement.

\begin{lemma}[The conductor-forty bound]\label{lem:conductor40}
If a multiset of $N$ roots of unity has sum $-a-b\sqrt{10}$, where
$a,b\in\mathbb Z_{\ge0}$, then $N\ge a+8b$. If equality holds and
the sum of the squared roots lies in $\mathbb Q(\sqrt{10})$, that
second sum is $a$.
\end{lemma}
\begin{proof}
Choose a cyclotomic field containing all roots and
$\mathbb Q(\sqrt{10})$. Take the union of the multiset over its Galois
group $G$ fixing $\sqrt{10}$, and put $M=|G|$. This gives a multiset
of $MN$ roots, with sum $-Ma-Mb\sqrt{10}$, whose multiplicities are
constant on $G$-orbits. Multiplicities remain integral throughout.

If an orbit consists of primitive $n$th roots and $40\nmid n$, its
sum is the Ramanujan sum $\mu(n)\in\{0,1,-1\}$ and its size is
$\varphi(n)$. Its rational contribution therefore costs at least one
root per unit. A negative contribution realizes equality only for
the singleton orbit $\{-1\}$.

If $40\mid n$, the primitive roots split into two equal orbits,
distinguished by the primitive quadratic character of conductor $40$.
Their sums are
\[
 \frac{\mu(n)\pm\tau_n(\chi_{40})}{2}.
\]
Here $\mu(n)=0$, and the induced Gauss sum vanishes unless
$n=40u$ with $u$ squarefree and coprime to $40$. In that case
\[
 \tau_n(\chi_{40})=\mu(u)\chi_{40}(u)\sqrt{40}.
\]
This formula follows by separating the factors coprime to $40$ in the
Chinese remainder theorem; an extra power of a prime dividing the
conductor gives zero. Thus an orbit contributes zero or
$\pm\sqrt{10}$, at cost at least eight roots. Equality for a nonzero
contribution requires $n=40$.

The rational and irrational contributions come from disjoint types
of orbit, so $MN\ge Ma+8Mb$. Equality leaves only $Ma$ copies of
$\{-1\}$ and $Mb$ copies of the negative quadratic orbit of primitive
fortieth roots. The latter orbit has squared sum zero: reduction of
its exponents modulo $20$ gives every unit once, since
$\chi_{40}(x+20)=-\chi_{40}(x)$, and the resulting sum is
$\mu(20)=0$. The averaged squared sum is therefore $Ma$.
By the additional hypothesis, Galois averaging multiplied the original
second sum by $M$. Dividing by $M$ gives the assertion.
\end{proof}

In our application $m,t$ are even and $a=m^2/2$, $b=mt/2$.
Lemma~\ref{lem:conductor40} gives
\[
 4m^2+2\ge\frac{m^2}{2}+4mt
\]
For $m\ge4$ one has $t/m>\sqrt{9/10}>15/16$, making the right side
strictly larger than $17m^2/4>4m^2+2$. The remaining $m=t=2$ is the
equality case,
so the squared sum must be $2$, whereas \eqref{eq:indicator-second}
gives $-2$ or $-2-4\sqrt{10}$. Thus $m=1$.
The same argument, including the indispensable second-moment
hypothesis, is recorded in
\artifact{hi_cft_20260923_next_structure_rank6_indicators.md} and
\artifact{hi_cft_20260923_next_rank6_root_audit.md}.

Consequently
\begin{equation}
 d=\frac{3+\sqrt{13}}2,\qquad
 s=\frac{13+3\sqrt{13}}2,\qquad
 d_X=\frac{11+3\sqrt{13}}2.
 \label{eq:haagerup-dimensions}
\end{equation}
\begin{lemma}[The six residual twists]\label{lem:six-residual-twists}
For $m=1$, there are six residual central simples, each with $n_j=1$.
Their twists are the six primitive thirteenth roots with quadratic
nonresidue exponents, and $\nu_2(\rho)=1$.
\end{lemma}
\begin{proof}
Equation~\eqref{eq:indicator-first} gives a multiset of six roots
with sum $(-1-\sqrt{13})/2$. Average it over a cyclotomic Galois group
fixing $\sqrt{13}$, as in Lemma~\ref{lem:conductor40}. An orbit with
nonzero $\sqrt{13}$ coefficient has coefficient of absolute value
$1/2$ and at least six members. To see the equality condition, the
induced quadratic Gauss-sum formula is nonzero only at primitive
orders $13u$, with $u$ squarefree and coprime to $13$. The orbit size
is $6\varphi(u)$, so equality permits only $u=1,2$.

The entire available length is already needed for the irrational
coefficient. Hence no rational-only orbit or cancellation can occur.
The two possible negative quadratic orbits, of orders $13$ and $26$,
have sums $(-1-\sqrt{13})/2$ and $(1-\sqrt{13})/2$, respectively.
The required rational coefficient forces every orbit to be the former.
Consequently every original root is one of the six nonresidue
thirteenth roots. These six roots are linearly independent over
$\mathbb Q$: a rational relation would give a polynomial of degree
at most twelve, with zero constant term, divisible by
$\Phi_{13}(t)=1+t+\cdots+t^{12}$, and hence zero.
Comparison with the quadratic period shows that each root occurs
once. As each object was repeated $n_j^2$ times, all six objects have
$n_j=1$ and distinct twists.

Multiplication by two sends nonresidues modulo thirteen to residues,
so the second sum is $(-1+\sqrt{13})/2$.
Equation~\eqref{eq:indicator-second} now gives $\varepsilon=1$.
\end{proof}

The center thus has twelve simples. Its dimensions and twist exponents are
\[
 \begin{array}{c|cccc}
 \text{number}&1&1&4&6\\\hline
 \text{dimension}&1&d^2&s&3d
 \end{array}
 \qquad
 0,0,0,0,\tfrac13,\tfrac23,
 \tfrac2{13},\tfrac5{13},\tfrac6{13},
 \tfrac7{13},\tfrac8{13},\tfrac{11}{13}.
\]
The complete $S,T$ data follow from
\cite[Theorem~A(2), equivalently Theorem~5.1]{GSY}: this statement assumes
a complex spherical untwisted HI category. The auxiliary metric group
has prime order $13$ and its negative quadratic Gauss sum fixes its
isometry class. This identifies the standard Haagerup modular data.
It does not by itself identify a braided category; that stronger
conclusion is treated separately below and in the analytic section.

\subsection{An optional algebraic classification route}

There is a second possible route from this fusion ring to an abstract
ribbon identification. We isolate its extra input rather than folding
an unreplayed classification into the preceding modular-data argument.
Vercleyen \cite[Section~4.4 and Table~9.1, p.~171]{Vercleyen} reports an
exhaustive classification of complex pivotal multiplicity-free fusion
categories of rank at most seven. Its list for $\operatorname{HI}(\mathbb
Z_3)$ has four entries. Gauge equivalence there permits complex nonzero
scalars and ring automorphisms; no unitarity assumption is imposed.
The solution method uses exact polynomial and algebraic-number
operations, with a separate numerical redundancy check. The complete
notebooks and decision logs are available on request rather than
included in the thesis \cite[Sections~4.1.7--4.1.8 and~4.4.1]{Vercleyen}.
We have not independently replayed that exhaustion calculation.

\begin{proposition}[Ribbon identification with an additional classification input]
\label{prop:optional-algebraic-identification}
Assume, in addition to the hypotheses used above, that the four-entry
list in \cite[Table~9.1]{Vercleyen} exhausts complex pivotal fusion
categories with the $\operatorname{HI}(\mathbb Z_3)$ fusion ring. Then
$\mathcal F$ is tensor equivalent to $\mathcal H_2$ or $\mathcal H_3$,
and there is an abstract ribbon equivalence
\[
 \operatorname{Rep}(U)\simeq Z(\mathcal H_3).
\]
In particular this modular category is abstractly unitarizable.
\end{proposition}
\begin{proof}
Evans--Gannon \cite[Theorem~3(ii)]{EG15} construct four inequivalent
pivotal categories with this ring. Two have noninvertible dimension
$(3+\sqrt{13})/2$ and are the unitary categories $\mathcal H_2$ and
$\mathcal H_3$; the other two have dimension $(3-\sqrt{13})/2$.
Their theorem classifies the endomorphism realizations considered
there; the additional exhaustion hypothesis is what makes this a
complete list for the present abstract $\mathcal F$.

A pivotal structure on any categorification of this ring is unique
when it exists. Indeed, the ratio of two such structures is a tensor
automorphism of the identity. Write its scalar on a simple $Z$ as
$t_Z$. Since $\rho^2$ contains both $\mathbf1$ and $\rho$,
$t_\rho^2=1=t_\rho$. It also contains $a\rho$ and $a^2\rho$,
so their scalars are one, and $a\otimes\rho=a\rho$ gives $t_a=1$.
Thus all six scalars are one. The negative-dimensional examples
cannot acquire a positive spherical structure by repivotalization.
The positive spherical category $\mathcal F$ is therefore one of the
two positive examples.

The categories $\mathcal H_2$ and $\mathcal H_3$ are Morita equivalent
by \cite[Theorem~1.4]{GS}, so their centers are braided equivalent.
Combining this with $\operatorname{Rep}(U)\simeq Z(\mathcal F)$ gives
the claimed braided equivalence. Both sides have the positive
spherical structures, which are unique; the induced pivotal and
twist structures are consequently preserved. This makes the
equivalence ribbon. Transporting the unitary structure from the
Haagerup center proves abstract unitarizability.
\end{proof}

This proposition does not establish positivity of the canonical VOA
fusion pairings or the analytic realization of the transported adjoint.
Those are separate complete-unitarity assertions addressed later.
Neither this classification input nor the proposition is needed for
the character computation that follows.

\section{The complete character vector}
\label{sec:characters}

Let $q=e^{2\pi i\tau}$. Character exponents are $h-c/24=h-1/3$.
Using the preceding twist list and strict positivity of nonvacuum
lowest weights, the only allowed nonvacuum pole is
$\gamma q^{-7/39}$ in the twist-$2/13$ coordinate. Here
$\gamma\in\mathbb Z_{\ge0}$ is its top dimension if that module starts
at $2/13$, and is zero if its lowest weight is shifted upward. The
vacuum pole is $q^{-1/3}$ with coefficient one.

\begin{lemma}[Uniqueness from the principal part]\label{lem:haagerup-character-unique}
Two character vectors for this modular representation with the same
$\gamma$ coincide in every weight.
\end{lemma}
\begin{proof}
The modular representation has finite image by the congruence theorem
\cite[Theorem~I]{DLN}. A constant vector must be fixed by both $S$ and
the \emph{character} $T$ matrix, which
includes $e^{-2\pi i/3}$. The only $T$-fixed coordinate is the object
of categorical twist $e^{2\pi i/3}$. Its $S$ image has a nonzero vacuum
component, so it cannot span an invariant line. There are therefore
no invariant constants, and Lemma~\ref{lem:pole-cancellation} applies.
\end{proof}

Evans--Gannon \cite[Sections~1.4 and~2.4, equation~(2.7)]{EG}
provide a fundamental matrix $\Xi$ whose
vacuum and $d_3$ columns have precisely the two principal parts above.
Thus, in their labeling,
\[
 \boldsymbol\chi=\Xi_{:,0}+\gamma\Xi_{:,d_3},
 \qquad\dim U_1=6+13\gamma.
\]
The actual six-current result forces $\gamma=0$. Their recursive
construction of $\Xi$ and the uniqueness lemma give the entire vector,
not just a finite coefficient match. In particular,
\begin{equation}
 \operatorname{ch}U
 =q^{-1/3}\bigl(1+6q+120q^2+956q^3+6010q^4+\cdots\bigr).
 \label{eq:vacuum-character}
\end{equation}
The lowest weights and top dimensions are
\[
\begin{array}{c|c|r@{\qquad}c|c|r}
 \text{label}&h&\dim M_{\rm top}&\text{label}&h&\dim M_{\rm top}\\\hline
 0&0&1&d_1&6/13&7\\
 b&1&80&d_2&11/13&42\\
 a&1&81&d_3&15/13&119\\
 c_0&1&81&d_4&5/13&5\\
 c_1&1/3&3&d_5&7/13&13\\
 c_2&2/3&27&d_6&8/13&14
\end{array}
\]
In particular the final maximum is $15/13$, improving the preliminary
$17/10$ bound. The proof record is
\artifact{hi_cft_20260923_next_structure_conditional_characters.md}.
Although this vector fixes every unflavoured homogeneous dimension,
it does not determine the separate branching functions for the
fourteen-module core $A_L$, and it was not used to prove properness.

\begin{corollary}[The four unmixed lattice restrictions]
\label{cor:unmixed-lattice-restrictions}
The two lattice modules labelled $(0,\pm1)\in(\mathbb Z/3)^2$
restrict to the simple $U$-module $c_1$, and the two labelled
$(\pm1,0)$ restrict to $c_2$.
\end{corollary}
\begin{proof}
The first pair has lowest weight $1/3$, top dimension three, and
weights in $1/3+\mathbb Z_{\ge0}$. The next possible summand top,
$4/3$, exceeds the established maximum $15/13$. The character list
contains a unique simple of lowest weight $1/3$, with top dimension
three, so semisimplicity gives precisely one copy of $c_1$.
The second pair has lowest weight $2/3$, top dimension $27$, and
next possible top $5/3$; the same argument gives $c_2$.
No additional finite-top irreducibility calculation is required.
\end{proof}
